\chapter*{要旨}
\addcontentsline{toc}{chapter}{要旨}

本プロジェクトは、廃プラスチックボトルの分別を対象として、ALOHA 双腕ロボットに多段階の長時間作業を学習させるものである。左作業腕が散在するボトルから対象を選択して把持し、ボトル本体を保持する。右作業腕はキャップを把持して回し、取り外す。その後、左腕はボトル本体を一方の回収箱へ、右腕はキャップを別の回収箱へ投入する。対象探索、双腕協調、精密接触、連続回転、物体分離、分別投入を連続して行うため、単純な一回把持より明らかに難度が高い。

初年度には、専用収集ソフトウェア、データ編集センター、データセットのバージョン管理、モデル学習、実機連続運用を結ぶ一連の工程を構築した。データ基盤には有効な ALOHA プロジェクトが 51 件、軌跡が 2,413 本、フレームが 2,907,804 枚あり、記録フレームレート換算で約 16.15 時間となる。実機運用モデルが実際に用いたのは、25 件の重複しないデータ保管単位、1,051 本の軌跡、879,852 フレーム、約 4.89 時間である。フィルタ後の学習対象は 844,102 フレーム、約 4.69 時間であった。\ev{E03,E04} 全数値検査では、状態・動作の非有限値、全要素がゼロの動作レコード、軌跡内の時刻逆転、フレーム番号の欠落は確認されなかった。ただし、すべての動画を全フレーム復号したわけではないため、「データに問題がない」とは断定しない。\ev{E05}

実機運用モデルは、上方全景、左手首、右手首の 3 視点と、双腕・グリッパの 14 次元状態を入力する。一回の推論で将来 50 制御周期分の双腕動作列を生成し、50 Hz で制御する。現在の動作列を実行している間に次の動作列を先行計算する。\ev{E06,E07} 保存モデルのパラメータ数は 838,358,468 である。学習履歴は 6,059 件、0--59,990 ステップ、約 51.4 時間に及び、学習誤差は 0.0677 から 0.000671 へ低下した。現場で使用したのは最終記録ではなく、ステップ 19,000 の保存モデルである。\ev{E08,E09} 検証集合、テスト集合、実機成功率は記録されていない。したがって、誤差低下は示教行動への近さを示すが、実機成功率の代替にはならない。

現場展示では、ボトル本体とキャップを別々の箱へ投入する一連の動作が観察された。また、1 時間を超える連続運転と、平均約 2 本/分の処理が観察された。\ev{E10} ただし、完全な動画、ボトル単位の計数表、固定試験手順、自動成功判定がないため、これらは\textbf{現場観測に基づく概算}であり、標準試験に基づくベンチマーク値や安定成功率ではない。明確な課題は、キャップのないボトルにも空回し動作を行うことと、一部のボトルを逆向きに持ち上げることである。監査では、キャップなし条件を明示した 6 件のデータ保管単位、158 軌跡のすべてで無条件の「キャップを回す」指示が使われていた。この不整合は空回しと整合する優先度の高いデータリスクであるが、唯一の原因とまでは断定できない。\ev{E11}

強化学習の探索では、洗浄タスク 835 軌跡、238,816 学習セグメント、連続 5 ラウンド、約 6,154 万パラメータの行動最適化モデルを確認した。オフライン検証の行動誤差は相対的に約 6.7\% 低下した。\ev{E15,E16} 同条件の基礎モデル比較と実機試験がないため、現段階で確認できるのは「学習閉ループを構築した」ことであり、「実機能力が向上した」ことではない。

次年度は、まず固定条件、ボトル単位の計数、段階判定、失敗分類を整備する。次に、「キャップなしでは回さない」という条件、ボトル口の向き、把持修正、接触段階のデータを追加し、基礎模倣学習モデルを強化する。その後、同一評価手順で強化学習の改善幅を検証し、最後に空ボトルを水管へ挿入するためのデジタルツイン、視覚位置決め、接触制御へ進む。ミリメートル級精度は現時点では設計目標であり、実測済み性能ではない。

\figfull[.97\textwidth]{evidence_grade.pdf}{成果のエビデンス区分。再確認可能な事実、現場観測に基づく概算、次段階で定量化すべき指標を分けて示す。学習誤差や一回の展示を安定成功率として扱わない。}

\section*{初年度成果の概要}
\begin{table}[H]
\centering
\small
\caption{確認できた成果と現時点の限界}
\begin{tabularx}{\textwidth}{>{\japanesesans\bfseries}p{30mm}p{41mm}X}
\toprule
項目 & 確認できた成果 & 現時点の限界\\
\midrule
一連のタスク & ボトル把持、開栓、ボトル/キャップ分別の連続動作を観察 & ボトル単位の計数と完全動画がなく、標準試験手順に基づく成功率は算出不可\\
連続運転 & 1 時間超、約 2 本/分との現場概算 & 信頼区間を伴う標準処理量ではない\\
データ資産 & 51 プロジェクト、2,413 軌跡、約 16.15 時間 & 実機モデルの重複なし学習部分は約 4.89 時間\\
実機モデル & 8.38 億パラメータ、3 視点、双腕連続行動 & 学習誤差に検証・テスト対照がない\\
工程反復 & 41 回の追跡可能な実行、14 構成 & 中断が多く、41 件の成功実験とは呼べない\\
解釈用記録 & アテンション重みの記録 8,223 件 & 結果ラベルと遮蔽実験がなく、因果説明には使えない\\
強化学習 & データ、保存モデル、オフライン検証の閉ループ & 実機で基礎モデルを上回ることは未証明\\
\bottomrule
\end{tabularx}
\end{table}
